\section{Convergence with Stochastic Gradients, Momentum, and Geometric Error}
\label{app:convergence}

The preceding geometry identifies an algorithm to which a Riemannian
convergence theorem can apply. This section separates exact quotient gradient
descent, biased stochastic updates, transported momentum, and arbitrary
adaptive directions. These are different algorithms. A proof for one is not
a proof for all four. In particular, a numerical observation that balance-only
matches an adaptive QuotientOpt implementation is not an exact theorem about
the momentum method in the proposal.

\subsection{Assumptions and a directly verifiable smoothness constant}
\begin{assumption}[Regular iterates and lower-bounded objective]
\label{ass:regular}
The iterates and retraction segments used in a stated theorem remain in a
compact regular neighborhood of $\mathcal K_{a,b}$. The objective
$f:\mathcal M_r\to\mathbb R$ is differentiable and bounded below by
$f_\star> -\infty$. A retraction radius $\rho>0$ is chosen so that all
steps in the theorem are well defined and have lift norm at most $\rho$.
\end{assumption}

\begin{assumption}[Retraction smoothness]
\label{ass:retraction-smooth}
For all used horizontal lifts $\xi$ and base points $X$,
\begin{equation}
 f(R_X(\xi))\leq f(X)+g_X(\operatorname{grad}_g f(X),\xi)
                   +\frac{L_R}{2}\|\xi\|_g^2.
 \label{eq:retraction-smooth}
\end{equation}
Here the norm of a quotient tangent equals the factor norm of its horizontal
lift. This is a bound on the actual retraction pullback, not merely a
geodesic-smoothness assertion on a different update.
\end{assumption}

\begin{lemma}[An ambient sufficient condition]
\label{lem:ambient-smoothness}
Suppose an ambient extension of $f$ has an $L_f$-Lipschitz gradient along
all relevant product segments, $\|\nabla f(X)\|\leq G_f$ at their base
points, and $\sigma_1(X)\leq b$. For the factor retraction in
Eq.~\eqref{eq:factor-retraction}, Assumption~\ref{ass:retraction-smooth}
holds for steps of lift norm at most $\rho$ with
\begin{equation}
 L_R=G_f+L_f\left(\sqrt{2b}+\frac\rho2\right)^2.
 \label{eq:explicit-smoothness}
\end{equation}
\end{lemma}
\begin{proof}
For a horizontal lift $\xi$, put $Z=\xi_BA+B\xi_A$ and
$Q=\xi_B\xi_A$. Equation~\eqref{eq:exact-retraction} gives
$R_X(\xi)-X=Z+Q$. Since
$\|Z\|\leq\sqrt{2b}\|\xi\|$ and $\|Q\|\leq\|\xi\|^2/2$,
ambient smoothness implies
\begin{align}
 f(R_X(\xi))&\leq f(X)+\langle\nabla f(X),Z\rangle
       +\langle\nabla f(X),Q\rangle+\tfrac{L_f}{2}\|Z+Q\|^2\\
 &\leq f(X)+g_X(\operatorname{grad}_g f(X),\xi)
       +\tfrac12\left[G_f+L_f(\sqrt{2b}+\rho/2)^2\right]\|\xi\|^2.
\end{align}
The first-order identity is Theorem~\ref{thm:metric-explicit}. This proves
the claim without replacing the factor retraction by an ambient Euler step.
\end{proof}

This lemma is useful for quadratic matrix losses, whose $L_f$ is controlled
by the measurement operator. For neural networks, bounded activations alone
do not automatically bound every derivative, particularly for PDE losses
containing higher derivatives. Their regularity assumptions must be checked
on the actual parameter neighborhood. For SIREN, the derivative bound
contains the frequency multiplier; it is not the same as the bound for
a one-Lipschitz activation.

\subsection{Deterministic descent and the PL regime}
\begin{theorem}[Deterministic first-order stationarity]
\label{thm:deterministic}
Under Assumptions~\ref{ass:regular}--\ref{ass:retraction-smooth}, let
$g_t=\operatorname{grad}_g f(X_t)$ and $X_{t+1}=R_{X_t}(-\eta_tg_t)$.
If $0<\eta_t\leq L_R^{-1}$, then
\begin{equation}
 f(X_{t+1})\leq f(X_t)-\frac{\eta_t}{2}\|g_t\|_g^2.
 \label{eq:deterministic-descent}
\end{equation}
Consequently
\begin{equation}
 \min_{t<T}\|g_t\|_g^2\leq
 \frac{2(f(X_0)-f_\star)}{\sum_{t<T}\eta_t}.
 \label{eq:deterministic-stationarity}
\end{equation}
For a constant admissible step this is $O(T^{-1})$ in squared gradient
norm and $O(T^{-1/2})$ in gradient norm.
\end{theorem}
\begin{proof}
Apply Eq.~\eqref{eq:retraction-smooth} with $\xi=-\eta_tg_t$.
The linear term is $-\eta_t\|g_t\|_g^2$, and the quadratic term is at
most $\eta_t\|g_t\|_g^2/2$. This proves
Eq.~\eqref{eq:deterministic-descent}. Sum it and use the lower bound
$f(X_T)\geq f_\star$; the minimum is at most the stepsize-weighted average.
\end{proof}

\begin{assumption}[Quotient Polyak--\L{}ojasiewicz inequality]
\label{ass:pl}
On the invariant neighborhood of interest, the same minimum $f_\star$
is attained and
\begin{equation}
 \tfrac12\|\operatorname{grad}_g f(X)\|_g^2
 \geq\mu_g(f(X)-f_\star),\qquad \mu_g>0.
 \label{eq:quotient-pl}
\end{equation}
\end{assumption}

\begin{theorem}[PL contraction and continuous-time dissipation]
\label{thm:pl}
Under the preceding assumptions and a constant admissible step,
\begin{equation}
 f(X_T)-f_\star\leq(1-\mu_g\eta)^T(f(X_0)-f_\star).
 \label{eq:pl-linear}
\end{equation}
The gradient flow $\dot X=-\mathcal K_X(\nabla f(X))$ satisfies
\begin{equation}
 \frac{d}{dt}f(X(t))=-\|\operatorname{grad}_g f(X(t))\|_g^2,
 \quad f(X(t))-f_\star\leq e^{-2\mu_gt}(f(X(0))-f_\star).
 \label{eq:flow-dissipation}
\end{equation}
If an embedded-metric PL inequality holds with constant $\mu_E$ on
the same domain, Eq.~\eqref{eq:gradient-equivalence} implies
$\mu_g\geq a\mu_E$. This comparison does not supply the missing PL
inequality for a general nonconvex network.
\end{theorem}
\begin{proof}
Insert Eq.~\eqref{eq:quotient-pl} into
Eq.~\eqref{eq:deterministic-descent} and iterate. Along the flow, the
chain rule and the definition of the Riemannian gradient give the
dissipation identity. Gronwall's inequality applied to
$\dot\Delta\leq-2\mu_g\Delta$ yields its exponential bound. Finally,
Eq.~\eqref{eq:gradient-equivalence} bounds the quotient squared gradient
below by $a$ times the embedded squared gradient.
\end{proof}

Balancing minimizes factor condition numbers on an orbit, whereas $a\mu_E$
depends on the represented product and the restricted loss landscape. These
are separate quantities. In particular, neither the PL constant nor an
improved iteration complexity follows from the factor condition lock alone.

\subsection{Stochastic stationarity with explicit bias}
Let $\mathcal F_t$ contain the iterate and optimizer history before drawing
the current stochastic sample. Write a horizontal estimator as
\begin{equation}
 \widehat g_t=g_t+b_t+\zeta_t,\qquad
 \mathbb E[\zeta_t\mid\mathcal F_t]=0,\qquad
 \mathbb E[\|\zeta_t\|_g^2\mid\mathcal F_t]\leq\nu_t^2.
 \label{eq:biased-estimator}
\end{equation}
The predictable vector $b_t$ is the conditional mean bias. It can include
the horizontal component of an inexact solve or a controlled coordinate
error; its size is not presumed zero because a program calls its output
a quotient gradient.

\begin{theorem}[Biased stochastic quotient descent]
\label{thm:stochastic}
Suppose Assumptions~\ref{ass:regular}--\ref{ass:retraction-smooth} hold
almost surely and $X_{t+1}=R_{X_t}(-\eta_t\widehat g_t)$ with
$0<\eta_t\leq(4L_R)^{-1}$. Then
\begin{align}
 \mathbb E[f(X_{t+1})\mid\mathcal F_t]
 &\leq f(X_t)-\frac{\eta_t}{2}\|g_t\|_g^2
     +2\eta_t\|b_t\|_g^2+\frac{L_R\eta_t^2}{2}\nu_t^2.
 \label{eq:biased-descent}
\end{align}
For $S_T=\sum_{t<T}\eta_t$,
\begin{equation}
 \frac{\sum_{t<T}\eta_t\mathbb E\|g_t\|_g^2}{S_T}
 \leq\frac{2\Delta_0}{S_T}
 +\frac{4\sum_{t<T}\eta_t\mathbb E\|b_t\|_g^2}{S_T}
 +\frac{L_R\sum_{t<T}\eta_t^2\nu_t^2}{S_T},
 \label{eq:stochastic-bound}
\end{equation}
where $\Delta_0=f(X_0)-f_\star$. This is also a bound for a random
output iterate chosen with probability $\eta_t/S_T$.
\end{theorem}
\begin{proof}
Conditional expectation in Eq.~\eqref{eq:retraction-smooth} gives
\begin{align}
 \mathbb E_t f(X_{t+1})
 &\leq f(X_t)-\eta_t\|g_t\|_g^2-\eta_t\langle g_t,b_t\rangle_g\\
 &\quad+\tfrac12L_R\eta_t^2\{\|g_t+b_t\|_g^2+\nu_t^2\}.
\end{align}
Use $|\langle g_t,b_t\rangle_g|\leq\|g_t\|_g^2/4+\|b_t\|_g^2$
and $\|g_t+b_t\|_g^2\leq2\|g_t\|_g^2+2\|b_t\|_g^2$.
The coefficient of $\eta_t\|g_t\|_g^2$ is at most
$-3/4+L_R\eta_t\leq-1/2$. The bias coefficient is at most
$1+L_R\eta_t\leq5/4\leq2$. This proves
Eq.~\eqref{eq:biased-descent}. Take total expectations, telescope the
objective terms, and divide by $S_T/2$ to obtain
Eq.~\eqref{eq:stochastic-bound}. The random-iterate interpretation is
the definition of its conditional sampling distribution.
\end{proof}

For constant $\nu_t\leq\nu$ and $\eta=c/\sqrt T$, the unbiased bound
is $2\Delta_0/(c\sqrt T)+L_Rc\nu^2/\sqrt T$, subject to the step-radius
restriction. A nonvanishing bias adds its nonvanishing mean-square floor.
To retain an $O(T^{-1/2})$ squared-stationarity rate, it is sufficient
that the weighted mean squared bias be $O(T^{-1/2})$; a fixed numerical
tolerance need not satisfy that condition at arbitrarily large horizons.

\begin{corollary}[Stochastic PL neighborhood]
\label{cor:stochastic-pl}
If Eq.~\eqref{eq:quotient-pl} also holds, $\|b_t\|_g\leq B$,
$\nu_t\leq\nu$, and $\eta$ is constant, then
\begin{equation}
 \mathbb E[f(X_T)-f_\star]
 \leq(1-\mu_g\eta)^T\Delta_0+
 \frac{2B^2+L_R\eta\nu^2/2}{\mu_g}.
 \label{eq:stochastic-pl-floor}
\end{equation}
\end{corollary}
\begin{proof}
Apply the PL inequality to the negative gradient term in
Eq.~\eqref{eq:biased-descent}, then sum the resulting affine recursion.
The geometric sum is bounded by $(\mu_g\eta)^{-1}$.
\end{proof}

\subsection{From geometric errors to the bias term}
Suppose a computation at a balanced point uses an approximate projected
direction. Proposition~\ref{prop:connection-error} bounds the factor error
by a quantity $E_{{\rm conn},t}$. Its horizontal projection has norm at
most $E_{{\rm conn},t}$ because the exact projector is nonexpansive.
However, the computed factor direction may also retain a vertical part.
Projecting that part away changes the implemented product by the exact
mixed-term expression in Eq.~\eqref{eq:correct-cross-terms}.
Consequently there are two errors to control: a tangent estimator error of
size $E_{{\rm conn},t}$ and a finite-step product error of size
$\eta_t^2q_t$. On a regular local chart with inverse differential norm
bounded by $C_{\rm chart}$, the latter can be represented as an additional
effective tangent error of order $C_{\rm chart}\eta_tq_t$, with higher
orders controlled by the chart radius.

The bound in Eq.~\eqref{eq:stochastic-bound} can therefore be used with
an explicitly justified estimate of the form
\begin{equation}
 \|b_t\|_g\leq E_{{\rm conn},t}+E_{{\rm coord},t}
       +E_{{\rm transport},t}+C_{\rm chart}\eta_tq_t.
 \label{eq:error-budget}
\end{equation}
This is an error accounting formula, not a claim that all implementations
have those errors bounded. In particular, computing gradients before a
nonorthogonal balance change and then treating them as gradients at the
new factors can give an uncontrolled $E_{{\rm coord},t}$. Correct
recomputation after balancing makes this specific coordinate error zero.
State transport without a uniform norm bound can likewise invalidate a
bounded-transport assumption near rank deficiency.

\subsection{Transported momentum: a complete tracking argument}
We next analyze the normalized momentum convention used in the proposal:
\begin{align}
 v_t&=\beta\mathcal T_{t-1\to t}v_{t-1}
          +(1-\beta)\widehat g_t,\qquad v_{-1}=0,\\
 X_{t+1}&=R_{X_t}(-\eta v_t),\qquad 0\leq\beta<1.
 \label{eq:momentum-algorithm}
\end{align}
For this theorem $\widehat g_t=g_t+\zeta_t$ is unbiased. The transport
is an isometry on the used tangents; parallel transport is one example.
Assume, on the same neighborhood,
\begin{equation}
 \|\mathcal T_{t-1\to t}g_{t-1}-g_t\|_g
       \leq L_Tc_R\eta\|v_{t-1}\|_g,
 \label{eq:transported-grad-lipschitz}
\end{equation}
and $\mathbb E_t\|\zeta_t\|_g^2\leq\nu^2$. These assumptions explicitly
include the gradient variation along the chosen update and transport.
An arbitrary coordinate transport through a large nonorthogonal gauge
change is not automatically an isometry.

\begin{lemma}[Momentum tracking recurrence]
\label{lem:momentum-tracking}
Let $e_t=v_t-g_t$, $q_t=\mathbb E\|e_t\|_g^2$, and
$V_t=\mathbb E\|v_t\|_g^2$. For $t\geq1$,
\begin{equation}
 q_t\leq\beta q_{t-1}+
 \frac{\beta^2L_T^2c_R^2\eta^2}{1-\beta}V_{t-1}
 +(1-\beta)^2\nu^2.
 \label{eq:tracking-recursion}
\end{equation}
Consequently
\begin{equation}
 \sum_{t<T}q_t\leq\frac{q_0}{1-\beta}+
 \frac{\beta^2L_T^2c_R^2\eta^2}{(1-\beta)^2}
       \sum_{t<T}V_t+T(1-\beta)\nu^2,
 \label{eq:tracking-sum}
\end{equation}
where $q_0\leq\beta^2\|g_0\|_g^2+(1-\beta)^2\nu^2$.
\end{lemma}
\begin{proof}
Set $d_t=\mathcal T_{t-1\to t}g_{t-1}-g_t$. Subtract $g_t$ from
Eq.~\eqref{eq:momentum-algorithm} to obtain
\begin{equation}
 e_t=\beta\mathcal T e_{t-1}+\beta d_t+(1-\beta)\zeta_t.
\end{equation}
The first two terms are measurable before the current sample. Their cross
term with $\zeta_t$ has zero conditional mean. Isometry and the weighted
Young inequality give
\begin{equation}
 \|\beta\mathcal T e_{t-1}+\beta d_t\|_g^2
 \leq\beta\|e_{t-1}\|_g^2+
       \frac{\beta^2}{1-\beta}\|d_t\|_g^2.
\end{equation}
Insert Eq.~\eqref{eq:transported-grad-lipschitz} to prove the recurrence.
Summing it from $1$ to $T-1$ and moving the $\beta\sum q_t$ terms to
the left proves Eq.~\eqref{eq:tracking-sum}; extending the nonnegative
$V_t$ sum to $T-1$ only enlarges its right side. At the first iteration,
$e_0=-\beta g_0+(1-\beta)\zeta_0$, whose expected squared norm gives
the initial bound.
\end{proof}

\begin{theorem}[Stationarity of transported horizontal momentum]
\label{thm:momentum}
Under the stated momentum assumptions, choose
\begin{equation}
 \eta\leq\min\left\{\frac1{2L_R},
 \frac{1-\beta}{2\beta L_Tc_R}\right\},
 \label{eq:momentum-step}
\end{equation}
where the second restriction is omitted for $\beta=0$. Then
\begin{equation}
 \frac1T\sum_{t<T}\mathbb E\|g_t\|_g^2
 \leq\frac{2\Delta_0}{\eta T}
       +\frac{q_0}{(1-\beta)T}+(1-\beta)\nu^2.
 \label{eq:momentum-rate}
\end{equation}
For deterministic gradients and fixed $\beta<1$, this yields an
$O(T^{-1})$ squared-stationarity bound. For stochastic gradients, taking
$1-\beta=T^{-1/2}$ and $\eta=cT^{-1/2}$ with a sufficiently small
constant $c$ gives an $O(T^{-1/2})$ bound for $T\geq4$.
\end{theorem}
\begin{proof}
The polarization identity gives
$2\langle g_t,v_t\rangle_g=\|g_t\|_g^2+\|v_t\|_g^2-\|e_t\|_g^2$.
Applying retraction smoothness to the step in
Eq.~\eqref{eq:momentum-algorithm} and using $L_R\eta\leq1/2$ yields
\begin{equation}
 \mathbb E f(X_{t+1})\leq\mathbb E f(X_t)
 -\tfrac\eta2\mathbb E\|g_t\|_g^2
 -\tfrac\eta4 V_t+\tfrac\eta2 q_t.
 \label{eq:momentum-descent}
\end{equation}
Sum this inequality and use Lemma~\ref{lem:momentum-tracking}.
The second stepsize restriction implies
$\beta^2L_T^2c_R^2\eta^2/(1-\beta)^2\leq1/4$.
The contribution of the resulting $V_t$ term is at most
$\eta\sum V_t/8$, which is absorbed by the negative
$-\eta\sum V_t/4$. Drop the remaining nonpositive term and divide by
$\eta T/2$ to prove Eq.~\eqref{eq:momentum-rate}.
For $1-\beta=T^{-1/2}$, the first two terms are $O(T^{-1/2})$ because
$q_0$ is bounded; the variance term is $\nu^2T^{-1/2}$.
\end{proof}

This is a conservative tracking theorem. For a fixed stochastic momentum
coefficient it displays a variance floor; it is not a claim that the
empirically used fixed $\beta=0.9$ must converge to zero gradient under
constant noise. Tighter fixed-momentum results require a different
Lyapunov argument and additional regularity. If the implemented transport
is only approximately isometric, its norm inflation and consistency
errors must be included in the tracking recurrence. A large transport
condition number cannot be dropped from that calculation.

\begin{corollary}[Momentum with a slightly norm-inflating transport]
\label{cor:approximate-momentum}
Replace isometry in Theorem~\ref{thm:momentum} by
$\|\mathcal T\xi\|\leq(1+c\eta)\|\xi\|$ and retain
Eq.~\eqref{eq:transported-grad-lipschitz}. If
\begin{equation}
 \beta(1+c\eta)^2\leq\frac{1+\beta}{2},\qquad
 \eta\leq\min\left\{\frac1{2L_R},
 \frac{1-\beta}{2\sqrt2\beta L_Tc_R}\right\},
 \label{eq:approximate-transport-step}
\end{equation}
then
\begin{equation}
 \frac1T\sum_{t<T}\mathbb E\|g_t\|_g^2
 \leq\frac{2\Delta_0}{\eta T}
 +\frac{2q_0}{(1-\beta)T}+2(1-\beta)\nu^2.
 \label{eq:approximate-momentum-rate}
\end{equation}
For $\beta=0$ the transport restrictions are unnecessary. The same
$O(T^{-1/2})$ stochastic schedule is available with a sufficiently small
constant when $c,L_T,c_R$ are uniform.
\end{corollary}
\begin{proof}
The weighted Young step in Lemma~\ref{lem:momentum-tracking} now gives
\begin{equation}
 q_t\leq\beta(1+c\eta)^2q_{t-1}
 +\frac{\beta^2L_T^2c_R^2\eta^2}{1-\beta}V_{t-1}
 +(1-\beta)^2\nu^2.
\end{equation}
The first coefficient is at most $\rho=(1+\beta)/2$ by assumption.
Summing therefore gives
\begin{equation}
 \sum_{t<T}q_t\leq\frac{2q_0}{1-\beta}
 +\frac{2\beta^2L_T^2c_R^2\eta^2}{(1-\beta)^2}\sum_{t<T}V_t
 +2T(1-\beta)\nu^2.
\end{equation}
The second restriction in Eq.~\eqref{eq:approximate-transport-step}
makes the coefficient of $\sum V_t$ at most $1/4$.
Insert this sum into Eq.~\eqref{eq:momentum-descent}; the velocity terms
are absorbed as in Theorem~\ref{thm:momentum}. Division by $\eta T/2$
proves Eq.~\eqref{eq:approximate-momentum-rate}.
\end{proof}

\begin{proposition}[Local applicability to balance--project transport]
\label{prop:implemented-transport}
Consider an exact horizontal factor step from a balanced pair in a compact
regular neighborhood. Copy the factor arrays as an ambient tangent at
the new unbalanced pair, change coordinates by a smooth near-identity
balancing gauge $Q$, and project horizontally at the new balanced pair.
If the step's horizontal velocity has norm at most $M$, this transport
has norm at most $1+C\eta^2M^2$ for sufficiently small steps, with
$C$ depending on the regular neighborhood. If the gradient is continuously
differentiable there, it also satisfies a bound of the form
Eq.~\eqref{eq:transported-grad-lipschitz}. Thus it fits
Corollary~\ref{cor:approximate-momentum} with local uniform constants.
The assertion concerns exact horizontal projection; damping and solve
errors require their additional terms in Eq.~\eqref{eq:error-budget}.
\end{proposition}
\begin{proof}
For a horizontal velocity $v$, the factors after the step have balance
defect
\begin{equation}
 (B-\eta v_B)^\top(B-\eta v_B)
 -(A-\eta v_A)(A-\eta v_A)^\top
 =\eta^2(v_B^\top v_B-v_Av_A^\top).
\end{equation}
The first-order term is zero because $C_p(v)=0$. The smooth local
balancing equation has an invertible symmetric-gauge derivative, a
Sylvester operator whose inverse is bounded on the compact regular
domain. The implicit-function bound therefore gives a balancing gauge
with $\|Q-I\|_2+\|Q^{-1}-I\|_2\leq C\eta^2M^2$.
The orthogonal ambiguity may be fixed by smooth Procrustes alignment;
its action is an isometry and does not change the norm estimate.
Equation~\eqref{eq:coordinate-transport} bounds the coordinate-change
operator norm by $1+C\eta^2M^2$, and horizontal projection has norm at
most one. This proves the transport bound.

At zero step the transport is the identity on horizontal tangents.
In a smooth frame, its difference from the identity is bounded by the
step displacement, and the continuously differentiable gradient has a
bounded derivative on the compact neighborhood. Applying the mean-value
bound to their composition yields a constant times
$\eta\|v\|$, which is Eq.~\eqref{eq:transported-grad-lipschitz}
after incorporating the local retraction displacement constant.
\end{proof}

The bounded-velocity premise is a substantive regularity condition,
particularly with stochastic gradients. For almost surely bounded oracle
gradients and a fixed finite horizon it can be checked from the momentum
recursion and the local transport bound; a horizon-independent bound
needs a small-step invariant velocity bound. It is not supplied by a
variance bound alone. These qualifications make the relationship between
the ideal isometric theorem and the implemented first-moment transport
explicit, instead of silently equating the two.

\subsection{Predictable preconditioners and the Adam boundary}
\begin{proposition}[Bounded predictable tangent preconditioning]
\label{prop:predictable-preconditioner}
Let $H_t$ be a self-adjoint tangent operator measurable with respect to
$\mathcal F_t$, with $\alpha I\preceq H_t\preceq\gamma I$, $\alpha>0$.
For the unbiased estimator $g_t+\zeta_t$, update
$X_{t+1}=R_{X_t}(-\eta H_t(g_t+\zeta_t))$.
If $\eta\leq\alpha/(L_R\gamma^2)$, then
\begin{equation}
 \frac1T\sum_{t<T}\mathbb E\|g_t\|_g^2
 \leq\frac{2\Delta_0}{\alpha\eta T}
            +\frac{L_R\eta\gamma^2\nu^2}{\alpha}.
 \label{eq:preconditioned-rate}
\end{equation}
\end{proposition}
\begin{proof}
Predictability makes the conditional mean update equal to $H_tg_t$.
The first-order descent is at least $\eta\alpha\|g_t\|_g^2$ and the
conditional squared update norm is at most
$\gamma^2(\|g_t\|_g^2+\nu^2)$. Retraction smoothness, the stepsize
restriction, and telescoping give the asserted result.
\end{proof}

An Adam denominator computed from the current sample is not predictable
in this sense, and its elementwise coordinates do not automatically define
a self-adjoint quotient operator with uniform spectral bounds. The present
proposition therefore does not prove convergence of ordinary AdamW,
balance-projected AdamW, or the archived adaptive heuristic. Those are
empirical comparators whose behavior must be measured. The momentum
algorithm has no elementwise second-moment state and has a different
justified analysis.

\subsection{Interpretation of long-horizon experiments}
The deterministic and stochastic bounds concern the minimum or a weighted
random iterate, not necessarily the last iterate. A best validation
checkpoint is selected by validation data and is not necessarily the
minimum training-gradient iterate. PL assumptions allow an objective-gap
statement for the last iterate, but they must be supplied separately.
Likewise, small raw gradient energy in the vertical space is an algebraic
identity and does not test convergence. We therefore log the quotient
gradient proxy, loss trajectory, numerical residuals, rank margin,
and actual task metric separately. For a near-zero $a$, the regular-domain
constants and the transport assumptions can fail even when a finite
training loss is reported.

The main practical consequence of the theory is a falsifiable set of
conditions: gradients must be evaluated in the current section, the
rank margin must be monitored, the joint connection residual must be small
for deep chains, and any claim of non-accumulation must be tied to the
actual update-map contraction. None of these conditions alone predicts
that balancing wins on every task family or at every learning rate.
